\documentclass[11pt]{article}
\usepackage{amsmath,amssymb,geometry,hyperref}
\geometry{margin=1in}
\newcommand{\Tr}{\operatorname{Tr}}
\begin{document}
\title{Exact fourth-order anticoherent states of fourteen symmetric qubits}
\author{Catalogue calculation}
\date{30 October 2026}
\maketitle

Let $|D_{14}^{r}\rangle$ be the normalized symmetric Dicke state with $r$
excitations, with $J_z|D_{14}^{r}\rangle=(7-r)|D_{14}^{r}\rangle$.
Put
\[
 A=\frac{7}{\sqrt{440}}+i\sqrt{\frac{39431}{320760}},\qquad
 B=\sqrt{\frac{364}{3645}},\qquad
 C=-\sqrt{\frac{91}{3000}}+i\sqrt{\frac{73229}{243000}}.
\]
Then the normalized state
\begin{align*}
 |\psi_+\rangle={}&A\bigl(|D_{14}^{1}\rangle+|D_{14}^{13}\rangle\bigr)\\
 &+B\bigl(|D_{14}^{4}\rangle+|D_{14}^{10}\rangle\bigr)
 +C|D_{14}^{7}\rangle
\end{align*}
is exactly fourth-order anticoherent. Its coefficientwise complex conjugate
$|\psi_-\rangle$ is another exact fourth-order anticoherent state.
These are two one-dimensional subspaces; their joint span is not claimed
to be anticoherent. No classification of all fourteen-qubit AC states is asserted.

\section{Exact verification}
The nonzero Dicke probabilities are
\[
(|c_1|^2,|c_4|^2,|c_7|^2,|c_{10}|^2,|c_{13}|^2)
=\frac1{3645}(854,364,1209,364,854).
\]
They sum to one. For any normalized Dicke coefficient vector $c$, its
$t$-qubit reduced state in the symmetric Dicke basis is $\rho_t=M_tM_t^\dagger$,
where
\[
 (M_t)_{ab}=c_{a+b}
 \sqrt{\frac{\binom ta\binom{14-t}b}{\binom{14}{a+b}}},
 \quad 0\le a\le t,\quad 0\le b\le14-t.
\]
Substituting the displayed probabilities makes every diagonal entry of
$\rho_4$ equal to $1/5$. Because the support has spacing three, its only
potentially nonzero off-diagonal entries are $(0,3),(1,4)$ and their conjugates.
Direct substitution gives
\[
 (\rho_4)_{03}=\frac B{1001}
 \left[\sqrt{286}(10A+\overline A)
       +\sqrt{42}(35\overline C+20C)\right],
 \qquad (\rho_4)_{14}=\overline{(\rho_4)_{03}}.
\]
Writing $A=a+ib$, the exact coefficient identity
\[
 C=\sqrt{\frac{143}{525}}\,(-a+3ib)
\]
makes the real and imaginary parts vanish separately. Thus
\[
 \boxed{\rho_4=\frac{I_5}{5}.}
\]
Complex conjugation preserves this equality, proving the claim for both signs.
The executable verifier independently forms every matrix entry with exact radicals.

These states have exactly order four, since the same computation gives
\[
 \Tr\rho_5^2=\frac{15209}{90750}
 =\frac16+\frac{14}{15125}>\frac16.
\]
They are therefore not fifth-order anticoherent.

\section{Derivation from the sparse numerical state}
On support $r=1,4,7,10,13$, the diagonal fourth-order AC conditions are
the five moment conditions through degree four. This is an invertible
Vandermonde system and yields the probabilities displayed above.
Use a global phase and a rotation about $z$ to make $c_1,c_4$ positive.
The remaining equations then admit the exact parametrization
\[
 z=\frac{-265\mathbin\pm27i\sqrt{39431}}{5368},\qquad |z|=1,
\]
\[
 c_1=\sqrt{\frac{854}{3645}},\quad
 c_4=\sqrt{\frac{364}{3645}},\quad
 c_7=\sqrt{\frac{17446}{273375}}(z-2),\quad
 c_{10}=z\sqrt{\frac{364}{3645}},\quad
 c_{13}=z^2\sqrt{\frac{854}{3645}}.
\]
A further $z$ rotation and global phase yield the palindromic formula above.
The supplied vector in \texttt{ACLiege.jl} agrees with the plus branch, after
such a rotation, to a maximum Dicke-amplitude error of approximately
$3.47\times10^{-19}$. This is a comparison with rounded numerical data,
separate from the exact existence proof. The number of printed decimal places
in the Julia file is not taken as an accuracy certificate.

\section{The two branches are distinct rotation orbits}
The support spacing gives a threefold rotation axis, and palindromicity
gives a perpendicular twofold rotation, so the proper rotation group contains
$D_3$. It is finite: invariance of the state ray under a continuous rotation
about an axis would make it an eigenvector of the corresponding spin component.
Its variance along that axis would vanish, contradicting second-order
anticoherence, which fixes every spin-component variance to $56/3$.

Its nonzero rank-five multipole excludes tetrahedral symmetry. Indeed the
spin-five representation has tetrahedral invariant multiplicity
\[
 \frac{\chi_5(0)+8\chi_5(2\pi/3)+3\chi_5(\pi)}{12}
 =\frac{11-8-3}{12}=0.
\]
Octahedral and icosahedral symmetry are excluded as well, since both contain
a tetrahedral subgroup. A larger dihedral group would require a rotation
of order greater than three about the existing threefold axis; the five
nonzero coefficients at spacing three prohibit it. Thus the full proper
rotation group is $D_3$, and its threefold axis is unique.

Any proper rotation relating the two branches must preserve this axis as an
unoriented line. Such rotations consist of $z$ rotations and their compositions
with a perpendicular half-turn. The quantity
\[
 R=\frac{c_1c_{13}}{c_4c_{10}}
\]
is invariant under their action and under a global phase. For the palindromic
state it equals $A^2/B^2$, whose imaginary part is nonzero. Conjugation replaces
$R$ by $\overline R$. Consequently the two conjugate branches are not related
by a collective proper rotation and a global phase.

\section{Maximum AC-subspace dimension}
For any fourth-order AC subspace of dimension $k$, apply Hilbert--Schmidt
orthonormal spin tensors of ranks zero, one and two to an orthonormal code basis.
Products of these tensors have rank at most four. The full AC compression
conditions therefore make the resulting $9k$ vectors orthogonal and nonzero.
They lie in the fourteen-qubit symmetric space of dimension fifteen, so
\[
 9k\le15,\qquad k\le1.
\]
Together with either exact state above this proves
\[
 \boxed{K(14,4)=1.}
\]
\end{document}
